% !PRINTED 386
\weberpage{386}

3.\ \gesperrt{If $\alpha_1, \alpha_2, \alpha_3, \alpha_4, \beta_1, \beta_2, \beta_3, \beta_4$ together denote all eight figures, then, in order to determine the influence of}
\[
\Pi_{\alpha_1\alpha_2\alpha_3\alpha_4} = \Pi_{\beta_1\beta_2\beta_3\beta_4}
\]
\gesperrt{on any root $[\mu\,\nu]$, one has to distinguish whether $\mu, \nu$ both occur among the $\alpha$ or both among the $\beta$, or whether the figure $\mu$ occurs among the $\alpha$ and $\nu$ among the $\beta$. In the first cases $[\mu\,\nu]$ passes into $[\mu'\,\nu']$, where $\mu, \nu, \mu', \nu'$ denote either all $\alpha$ or all $\beta$; in the third case $[\mu\,\nu]$ remains unchanged.}

In all cases the permutations $\Pi_{\alpha_1\alpha_2\alpha_3\alpha_4}$ are only of the second degree, and hence on a single repetition lead back to the identity.

The group $P$ is thereby completely determined and represented, and to establish the laws of composition in $P$ only a few more formulae are needed, which follow easily from the rule set up above. Here it is to be remarked that one can always take two distinct permutations $\Pi_{\alpha_1\alpha_2\alpha_3\alpha_4}$ in such a way that they have two or three figures in common in the index, since, if they have only one figure in common, one may take for the index of one of them its complement.

Now let
\[
\sigma = \begin{pmatrix} 1, & 2, & 3, & 4, & 5, & 6, & 7, & 8 \\ \alpha_1, & \alpha_2, & \alpha_3, & \alpha_4, & \alpha_5, & \alpha_6, & \alpha_7, & \alpha_8 \end{pmatrix}
\]
be any permutation of the eight figures, and let $(\alpha, \beta)$ denote the transposition of the two figures $\alpha$ and $\beta$, hence an element of $S$. Then
\weq{8}{\Pi_{1,2,3,4}\,\sigma = \sigma\,\Pi_{\alpha_1\alpha_2\alpha_3\alpha_4},}
\weq{9}{\Pi_{1,2,3,4}\,\Pi_{1,2,3,4} = 1,}
\weq{10}{\Pi_{1,2,3,4}\,\Pi_{1,2,3,5} = (4,5)\,\Pi_{1,2,3,4},}
\weq{11}{\Pi_{1,2,3,4}\,\Pi_{1,2,5,6} = (1,2)(3,4)(5,6)(7,8)\,\Pi_{1,2,7,8}.}

One proves these formulae easily by rule 3., by going through the individual cases, in which of course only quite a small number of types have to be considered; thus, for example, $[1\,4]$ passes by $\Pi_{1,2,3,4}$ into $[2\,3]$, and this by $\Pi_{1,2,3,5}$ into $[1\,5]$, and $[1\,5]$ is no longer altered by $\Pi_{1,2,3,4}$; consequently there is effected by
\weq{12}{\Pi_{1,2,3,4}\,\Pi_{1,2,3,5}\,\Pi_{1,2,3,4}}
%CONT>>
